\documentclass[11pt,letterpaper]{article}
\PassOptionsToPackage{hyphens}{url}
\usepackage[margin=.83in,top=.8in,bottom=.8in,headheight=15pt]{geometry}
\usepackage[T1]{fontenc}
\usepackage{lmodern,amsmath,amssymb,mathtools,bm,booktabs,array,graphicx,tikz}
\usetikzlibrary{arrows.meta,positioning,calc,decorations.pathreplacing,patterns}
\definecolor{ink}{HTML}{17344D}\definecolor{accent}{HTML}{087E8B}\definecolor{pale}{HTML}{EDF5F5}\definecolor{gold}{HTML}{AD6B19}
\providecommand{\dd}{\,\mathrm{d}}\providecommand{\R}{\mathbb{R}}
\providecommand{\E}{\mathbb{E}}\providecommand{\Var}{\operatorname{Var}}
\providecommand{\Cov}{\operatorname{Cov}}\providecommand{\diag}{\operatorname{diag}}
\providecommand{\norm}[1]{\left\lVert#1\right\rVert}
\providecommand{\abs}[1]{\left|#1\right|}
\usepackage{microtype,hyperref,fancyhdr,enumitem,amsthm,sectsty,needspace}
\allsectionsfont{\raggedright}
\hypersetup{colorlinks=true,linkcolor=ink,urlcolor=accent,pdfauthor={Zan Ahmad},pdftitle={Muscle force and the pressure--volume cycle / Review}}
\setlength{\parindent}{0pt}\setlength{\parskip}{6pt}\linespread{1.07}
\setlist{itemsep=3pt,topsep=4pt,leftmargin=*}
\emergencystretch=2em\clubpenalty=10000\widowpenalty=10000
\pagestyle{fancy}\fancyhf{}\fancyhead[L]{\small Mathematical Cardiology}\fancyhead[R]{\small Lecture 3}
\fancyfoot[L]{\footnotesize Zan Ahmad}\fancyfoot[R]{\footnotesize\thepage}\renewcommand{\headrulewidth}{.3pt}
\newtheorem{proposition}{Proposition}\newtheorem{theorem}{Theorem}\newtheorem{lemma}{Lemma}
\theoremstyle{definition}\newtheorem{definition}{Definition}\newtheorem{example}{Example}
\numberwithin{equation}{section}
\begin{document}
\section*{Muscle force and the pressure--volume cycle --- review sheet}
Use this after attempting the worked examples and practice. Explain each claim in words before calculating.

\section*{Key ideas}
\begin{itemize}
\item Filaments slide; they do not need to shorten themselves.
\item Force can be nonzero without shortening.
\item Valve states and conservation explain the four loop phases.
\item Stroke volume is a difference; ejection fraction is a ratio; output also includes rate.
\item Loop area represents net pressure-volume work, not total metabolic energy.
\end{itemize}
\section*{Common errors and repairs}
\textbf{A vertical loop segment means no muscle activity.} It means no chamber-volume change in this idealization; pressure and active force can change.

\textbf{A larger ejection fraction is an energy-efficiency measurement.} It is a volume fraction, not a ratio of mechanical to chemical energy.

\section*{Retrieve without your notes}
\begin{itemize}
\item Which axis carries volume on a pressure-volume loop?
\item What evidence would test a prescribed loop shape?
\end{itemize}
\section*{Readiness check}
\begin{itemize}
\item Describe sliding filaments without saying the filaments themselves shrink.
\item Identify filling, pressure-rise, ejection and pressure-fall phases from a loop.
\item Calculate stroke volume, ejection fraction and output with units.
\item Estimate the net work of a rectangular loop and distinguish it from total metabolic energy.
\end{itemize}
\end{document}
